
\documentclass[a4paper,12pt]{article}

\usepackage{JYM-harj}
\usepackage{tikz}
\usepackage{amsmath}

\setlength{\parindent}{0in}

\setlength{\voffset}{-1.0cm}
\addtolength{\textheight}{2.0cm}

\setlength{\hoffset}{-1.0cm}
\addtolength{\textwidth}{2.0cm}

\usepackage{graphicx}
\usepackage{verbatim}

\usepackage{hyperref}

%\usepackage[finnish]{babel}
%\usepackage[utf8]{inputenc}
%\usepackage[T1]{fontenc}
%\usepackage{amsmath,amsfonts,amssymb,amsthm,amscd}


% 1: Teht\"av\"at
% 2: Kaikki ratkaisut
% 3: T\"ahdett\"om\"at ratkaisut
% 4: T\"ahdelliset ratkaisut
\newcommand{\tversio}{1}

\newcommand{\harjoitusnro}{3.3.2026 klo 17:00}

\ifthenelse{\tversio = 1}{\pagestyle{empty}}{}

\begin{document} 

\otsikko

\begin{enumerate}

\item 
    \begin{enumerate}
    \end{enumerate}
    
\item 
\begin{enumerate}
\end{enumerate}

\item

\item Ratkaise yhtälö $z^7=1$ kompleksilukujen joukossa. Anna ratkaisut eksponentti- sekä napakoordinaattimuodossa. (6p)

\textbf{Vastaus:} Yhtälö voidaan esittää muodossa $r^7e^{7\phi i}=e^{0i}$. Muodostetaan yhtälöpari:

\begin{equation*}
    \begin{cases}
    r^7=1 \\
    7\phi = k \cdot 2\pi 
    \end{cases}
    \Longleftrightarrow
    \begin{cases}
        r=\sqrt[7]{1}=1 \\
        \phi=k\cdot \frac{2\pi}{7}
    \end{cases}
\end{equation*}
Näistä saadaan vastaukset:
\begin{enumerate}
    \item \textbf{Eksponenttimuodossa:} $e^{0i}$, $e^{\frac{2\pi}{7}i}$, $e^{\frac{4\pi}{7}i}$, $e^{\frac{6\pi}{7}i}$, $e^{\frac{8\pi}{7}i}$, $e^{\frac{10\pi}{7}i}$, $e^{\frac{12\pi}{7}i}$ \\
    \item \textbf{Napakoordinaattimuodossa:} $\cos(0)+i \sin(0)$, $\cos(\frac{2\pi}{7})+i \sin(\frac{2\pi}{7})$, $\cos(\frac{4\pi}{7})+i \sin(\frac{4\pi}{7})$, $\cos(\frac{6\pi}{7})+i \sin(\frac{6\pi}{7})$, $\cos(\frac{8\pi}{7})+i \sin(\frac{8\pi}{7})$, $\cos(\frac{10\pi}{7})+i \sin(\frac{10\pi}{7})$, $\cos(\frac{12\pi}{7})+i \sin(\frac{12\pi}{7})$
\end{enumerate}

\end{enumerate}

\end{document}